\scanpage{13}
\mspage{6}
p.ex. que la dim coh. de $\pi_i$ est égale à $n$, celle de $E_i$ est $\leq n+2$ (pour des coeff. \deleted{$\neq2$} de $m$-torsion, quand $m$ est impair) --- et par passage à la limite, des majorations correspondantes pour les dim coh. de $\pi_{\bar K/K}$ et de $E_{\bar K/K}^{\unsure{\ell}}$
\[
(14)\qquad \dim\operatorname{coh}\pi_{\bar K/K}\leq n,\qquad
\dim\operatorname{coh}_{\ell}E_{\bar K/K}\leq n+2\quad\text{si }\ell\neq2
\qquad(\ell\text{ nb premier})
\]
qui sont en fait même des \uline{égalités} (sauf erreur), et donnent donc une description cohomologique simple du degré de transcendance absolu de $K$.

\uline{Théorème 1} Soit $K$ un corps \inserted{extension} de type fini de $\mathbb Q$. \deleted{Alors} $\bar K$ une clôture algébrique de $K$. Alors pour tout sous-groupe ouvert $\unsure{\mathcal E|\epsilon}$ de $E_{\bar K/K}$, son centralisateur dans $E_{\bar K/K}$ est réduit au groupe unité. Idem pour $\pi_{\bar K/K}$.

Démonstration. Soit \deleted{$\unsure{E'\subset E}$}\inserted{$\Gamma'\subset\Gamma\subset$}$\Gamma_{\overline{\mathbb Q}/\mathbb Q}$ l'image de $\unsure{E|\mathcal E}$ dans $\Gamma$\inserted{$=\Gamma_{\bar K/K}$}, qui est donc un ss-groupe ouvert. \inserted{L'image dans $\Gamma$ du centralisateur de $\unsure{\mathcal E|\epsilon}$ dans $\unsure{\mathcal E|E}$ est contenue dans le centralisateur de $\Gamma'$ dans $\Gamma$.} \deleted{Si $g\in\Gamma_{\bar K/K}$ centralise $\unsure{\mathcal E|\epsilon}$, son image $\gamma$ dans $\Gamma$ centralise $\Gamma'\subset\Gamma_{\mathbb Q}$} Je dis qu'il est
\margin{12}
